% sec_ground_truth.tex

\subsection{Ground Truth Robustness}
\label{sec:ground_truth}

Is the FL--CADE association driven by market proximity alone, or do FL
firms actually behave differently from other always-losers? Three pieces
of evidence point to the latter. In CADE-present tenders, FL losing
bids have a mean log-spread of 3.08 above the winner, compared with
1.84 for non-FL bids ($t = 12.1$, $p < 0.001$)---even within the same
cartel-present tenders, FL firms bid from a different distribution.
Once we condition on participation count, FL firms' share of
convicted-firm tenders becomes indistinguishable from that of non-FL
always-losers ($p = 0.93$), confirming that the unconditional FL--CADE
link reflects high participation volume; the permutation test
(3.5$\times$, Section~\ref{sec:cade}) already accounts for this.
Finally, CADE overlap concentrates \emph{entirely} among FL firms:
non-FL always-losers with cover-bid-like spreads show zero CADE
overlap, while FL firms \emph{without} such spreads still show 6.2\%
overlap. It is the participation anomaly, not the bid pattern, that
drives the CADE association.
